\begin{definitionbox}{Definition (Presburger Signature)}
\begin{definition}[Presburger Signature]
\label{def:presburger-signature}
The \emph{Presburger signature} is the first-order signature
\[
\mathcal{L}_{\mathrm{Pres}} = \{0, S, +, <\},
\]
where $0$ is a constant symbol, $S$ is a unary function symbol, $+$ is a
binary function symbol, and $<$ is a binary relation symbol. Equality is the
logical equality symbol of first-order logic, not an additional nonlogical
relation symbol of the signature.
\end{definition}
\end{definitionbox}

\begin{remark*}[Standard quantified statement]
\[
\operatorname{PresburgerSignature}(\mathcal{L})
\Longleftrightarrow
\mathcal{L} = \{0,S,+,<\}
\land \operatorname{ar}(S)=1
\land \operatorname{ar}(+)=2
\land \operatorname{ar}(<)=2.
\]
\end{remark*}

\begin{remark*}[Interpretation]
The Presburger signature names addition and order, but it deliberately omits
multiplication as a function symbol. This makes Presburger arithmetic a theory
of the additive structure of the natural numbers rather than a theory of full
arithmetic.
\end{remark*}

\begin{dependencies}
\begin{itemize}
  \item \hyperref[def:signature]{First-Order Signature}
  \item \hyperref[def:arity]{Arity}
  \item \hyperref[def:relation-symbol]{Relation Symbol}
  \item \hyperref[def:signature-constant-symbol]{Constant Symbol}
  \item \hyperref[def:signature-function-symbol]{Function Symbol}
\end{itemize}
\end{dependencies}

\begin{definitionbox}{Definition (Presburger Arithmetic)}
\begin{definition}[Presburger Arithmetic]
\label{def:presburger-arithmetic}
\emph{Presburger arithmetic} is the complete first-order theory of the
standard natural-number structure in the Presburger signature. Let
\[
\mathcal N_{\mathrm{Pres}}
:=
(\mathbb{N},0,S,+,<).
\]
Then
\[
\mathrm{PA}_{\mathrm{Pres}}
:=
\operatorname{Th}(\mathcal N_{\mathrm{Pres}})
=
\left\{
\varphi:
\varphi \text{ is an }\mathcal L_{\mathrm{Pres}}\text{-sentence and }
\mathcal N_{\mathrm{Pres}}\models\varphi
\right\}.
\]
\end{definition}
\end{definitionbox}

\begin{remark*}[Standard quantified statement]
\[
\operatorname{PresburgerArithmetic}(T)
\Longleftrightarrow
T=
\left\{
\varphi:
\varphi \text{ is an }\mathcal L_{\mathrm{Pres}}\text{-sentence and }
(\mathbb{N},0,S,+,<)\models\varphi
\right\}.
\]
\end{remark*}

\begin{remark*}[Interpretation]
Presburger arithmetic isolates the part of arithmetic governed by addition.
It is strong enough to express linear equations, congruences, and order
properties of natural numbers, but it cannot directly express multiplication
as a primitive operation. The later Peano-system material supplies the more
general recursive framework from which addition, multiplication, order, and
induction are developed structurally.
\end{remark*}

\begin{dependencies}
\begin{itemize}
  \item \hyperref[def:presburger-signature]{Presburger Signature}
  \item \hyperref[def:natural-numbers]{Natural Numbers}
  \item \hyperref[def:addition-on-peano-system]{Addition on a Peano System}
  \item \hyperref[def:peano-system]{Peano System}
\end{itemize}
\end{dependencies}
